%% =============================================================================
%%  Etapa 3 — Estrutura
%% =============================================================================
\chapter{Estrutura}\label{cap:estrutura}

\begin{objetivo}
Descobrir como o trecho está construído: suas partes, as correspondências
entre elas e as formas literárias que usa. A forma faz parte da mensagem.
\end{objetivo}

\section{Paralelismo antitético}

Os vv.~20–26 formam dois blocos de quatro, espelhados. Cada felicidade tem um
\enquote{ai} correspondente, com a situação invertida; a conjunção
\gr{πλήν} (\enquote{porém}) faz de dobradiça entre eles.

\begin{center}
\begin{tikzpicture}[
    x=1mm, y=1mm,
    caixa/.style={draw=#1, fill=#1!6, rounded corners=2pt, text width=44mm,
      inner sep=4pt, align=left, font=\small, minimum height=15mm},
    ref/.style={font=\EBfnumeros\bfseries\footnotesize, text=rubrica},
    liga/.style={{Stealth[length=4pt]}-{Stealth[length=4pt]}, draw=destaque2, thick},
    rot/.style={font=\sffamily\scriptsize, text=suave, fill=white, inner sep=1.5pt,
      text width=20mm, align=center}]
  \node[font=\EBfversal\small, text=grego] at (-39, 17) {Felicidades};
  \node[font=\EBfversal\small, text=rubrica] at (39, 17) {Ais};
  \foreach \n/\y/\a/\b/\ra/\rb/\c in {
      1/0/{\gr{οἱ πτωχοί} — o Reino é vosso}/{\gr{οἱ πλούσιοι} — já tendes a consolação}/20b/24/{pobres\\×\\ricos},
      2/-20/{\gr{οἱ πεινῶντες νῦν} — sereis saciados}/{\gr{οἱ ἐμπεπλησμένοι νῦν} — tereis fome}/21a/25a/{fome\\×\\fartura},
      3/-40/{\gr{οἱ κλαίοντες νῦν} — rireis}/{\gr{οἱ γελῶντες νῦν} — chorareis}/21b/25b/{choro\\×\\riso},
      4/-64/{odiados, excluídos, insultados — \enquote{assim faziam aos \emph{profetas}}}/{elogiados por todos — \enquote{assim faziam aos \emph{falsos profetas}}}/22–23/26/{rejeição\\×\\aprovação}}{
    \node[caixa=grego] (f\n) at (-39, \y) {\a};
    \node[caixa=rubrica] (a\n) at (39, \y) {\b};
    \node[ref, anchor=south east] at (f\n.north east) {\ra};
    \node[ref, anchor=south east] at (a\n.north east) {\rb};
    \draw[liga] (f\n) -- node[rot]{\c} (a\n);
  }
  \node[draw=destaque2, fill=fundo, circle, inner sep=1pt, font=\EBfgrego\small,
    text=tinta] at (0, 11) {\gr{πλήν}};
\end{tikzpicture}
\end{center}

\paragraph{Observações.}
Os três primeiros pares são curtos e simétricos, com \gr{νῦν} nos
pares~2 e~3. O quarto par é mais longo e traz o fundamento histórico: o modo
como \enquote{os pais deles} tratavam os profetas. A última felicidade termina
com uma ordem (\enquote{alegrai-vos}); o último ai, com uma comparação que
desmascara a aprovação geral.

\begin{pergunta}
Onde está o \enquote{centro} do trecho? No par mais longo (22–23 × 26) ou
na dobradiça \gr{πλήν}? Que diferença isso faz para a leitura?
\end{pergunta}

\section{Formas literárias}

O trecho combina duas formas conhecidas da Bíblia hebraica e do judaísmo:
uma vem da sabedoria, a outra da profecia.

\subsection{O macarismo}
A forma \enquote{feliz quem\ldots, porque\ldots} é comum na literatura
sapiencial (\rb{Sl 1:1; 32:1–2; Pr 3:13}). Ali ela descreve quem vive com
sabedoria. Nos macarismos de Jesus, a razão (\gr{ὅτι}) está no \emph{futuro}
e depende de Deus: são macarismos \emph{escatológicos}, que anunciam uma
inversão, e não conselhos de conduta.

\subsection{O \enquote{ai} profético}
Os ais retomam uma forma dos profetas, que os usavam em séries contra a
injustiça e a autossuficiência: \rb{Is 5:8–23} (seis ais contra quem acumula
casas e campos), \rb{Am 6:1} (\enquote{ai dos que vivem tranquilos em Sião}),
\rb{Hc 2:6–19} (cinco ais contra o opressor). Ao unir macarismos e ais, Jesus
junta a voz do sábio e a voz do profeta.

\begin{nota}[Bênção e maldição]
O par felicidade/ai lembra a estrutura de bênçãos e maldições da aliança
(\rb{Dt 27–28}) e o \enquote{caminho dos justos} e o \enquote{caminho dos
ímpios} de \rb{Sl 1}. Mas note: os ais de Lucas não são maldições
(\gr{ἐπικατάρατος}); são lamentos e advertências.
\end{nota}

\section{Marcadores de discurso}

\begin{ebtabela}{@{}l l >{\raggedright\arraybackslash}X@{}}
\EBcab{Marcador} & \EBcab{Onde} & \EBcab{Função}\\\midrule
\gr{καὶ αὐτός} & 20 & introduz a fala; retoma Jesus como sujeito após a cena de 6:17–19\\
\gr{ὅτι} (6×) & 20–25 & dá a razão de cada felicidade e de cada ai\\
\gr{νῦν} (4×) & 21, 25 & contrasta o presente com o futuro\\
\gr{ὅταν} (2×) & 22, 26 & introduz situações eventuais nos pares longos\\
\gr{ἰδοὺ γάρ} & 23 & chama a atenção para a recompensa\\
\gr{κατὰ τὰ αὐτὰ γάρ} (2×) & 23, 26 & fundamenta com o passado (profetas / falsos profetas)\\
\gr{πλήν} & 24 & dobradiça adversativa\\
\end{ebtabela}

\begin{pergunta}
A mesma frase (\gr{κατὰ τὰ αὐτὰ γὰρ ἐποίουν \ldots\ οἱ πατέρες αὐτῶν})
fecha as duas metades. O que ela diz sobre o critério usado para julgar
\enquote{bem} ou \enquote{mal} alguém que fala em nome de Deus?
\end{pergunta}
